% Part: first-order-logic
% Chapter: introduction
% Section: models-theories

\documentclass[../../../include/open-logic-section]{subfiles}

\begin{document}

\olfileid{fol}{int}{mod}

\olsection{નિદર્શો અને સિદ્ધાંતો}

પ્રથમ-ક્રમ તર્કશાસ્ત્રની વાક્યરચના અને અર્થવિચાર વ્યાખ્યાયિત કરી લઈએ
પછી !!{structure}sના અને અર્થવિચારી ખ્યાલોના ગુણધર્મો તપાસવાનું શરૂ
કરી શકીએ છીએ. !!{derivation} તંત્રો પણ વ્યાખ્યાયિત કરીને તેમની તપાસ
કરી શકીએ. !!{sentence}sના કોઈ ગણ માટે પૂછી શકીએ: કઈ
!!{structure}s એ ગણમાંનાં બધાં !!{sentence}sને સત્ય બનાવે છે?
!!{sentence}sનો ગણ~$\Gamma$ આપેલો હોય, ત્યારે તેમને સંતોષતી
!!a{structure}~$\Struct{M}$ને \emph{~$\Gamma$નો નિદર્શ} કહે છે. આપણે
~$\Gamma$થી શરૂ કરીને તેના નિદર્શો શોધવાનો પ્રયત્ન કરી શકીએ---તેઓ
કેવા દેખાય છે? કેટલા મોટા કે નાના હોવા જરૂરી છે? પરંતુ કોઈ એક
!!{structure} અથવા !!{structure}sના કોઈ સંગ્રહથી શરૂ કરીને પણ પૂછી
શકીએ: તેમાં કયાં !!{sentence}s સત્ય છે? શું એવાં !!{sentence}s છે જે
આ !!{structure}sનું એ અર્થમાં \emph{લક્ષણનિર્ધારણ} કરે કે તે બરાબર
એ જ સંરચનાઓમાં સત્ય હોય? આવા પ્રશ્નો \emph{નિદર્શસિદ્ધાંત}ના ક્ષેત્રમાં
આવે છે. આ પ્રશ્નો \emph{સ્વયંસિદ્ધિમૂલક પદ્ધતિ}નો પણ પાયો છે: કોઈ
સિદ્ધાંતની સ્વયંસિદ્ધિઓ એટલે !!{sentence}sના ગણ વડે
!!{structure}sના સંગ્રહનું વર્ણન કરવું. સ્વયંસિદ્ધિઓમાંથી પ્રથમ-ક્રમ
તર્કશાસ્ત્રમાં ફલિત થતાં બરાબર બધાં !!{sentence}s સ્વયંસિદ્ધિઓના
દરેક નિદર્શમાં સત્ય છે, એ નિરીક્ષણ આ પદ્ધતિ શક્ય બનાવે છે.

એક બહુ સાદા દાખલા તરીકે પૂર્વક્રમો વિચારો. કોઈ ગણ~$A$ પરનો સંબંધ~$R$
સ્વવાચક અને પરંપરિત બંને હોય તો તે પૂર્વક્રમ છે. ગણ~$A$ અને તેના
પરનો દ્વિસ્થાની સંબંધ $R \subseteq A \times A$ મળીને એક દ્વિસ્થાની
સંબંધસંકેત~$\Obj P$ ધરાવતી પ્રથમ-ક્રમ ભાષા માટે !!a{structure} આપવા
જોઈતું બરાબર બધું આપે છે: આપણે $\Domain{M} = A$ અને
$\Assign{\Obj P}{M} = R$ રાખીએ. $R$ પૂર્વક્રમ હોવાથી તે સ્વવાચક અને
પરંપરિત છે, અને આ કહેતાં પ્રથમ-ક્રમ તર્કશાસ્ત્રનાં !!{sentence}sનો
ગણ~$\Gamma$ આપણે શોધી શકીએ છીએ:
\begin{align*}
  & \lforall[\Obj v_0][\Atom{\Obj P}{\Obj v_0,\Obj v_0}]\\
  & \lforall[\Obj v_0][\lforall[\Obj v_1][\lforall[\Obj v_2][((\Atom{\Obj P}{\Obj v_0,\Obj v_1} \land \Atom{\Obj P}{\Obj v_1,\Obj v_2}) \lif \Atom{\Obj P}{\Obj v_0,\Obj v_2})]]]
\end{align*}
આ !!{sentence}s માત્ર ``કોઈ પણ~$x$ માટે $Rxx$'' ($R$ સ્વવાચક છે) અને
``જ્યારે પણ $Rxy$ અને $Ryz$, ત્યારે $Rxz$ પણ'' ($R$ પરંપરિત છે)ના
સંકેતીકરણો છે. આપણે જોઈએ છીએ કે !!a{structure}~$\Struct{M}$ આ બે
!!{sentence}sના ગણ~$\Gamma$નો નિદર્શ ત્યારે અને માત્ર ત્યારે છે
જ્યારે $R$ (એટલે કે~$\Assign{\Obj P}{M}$),~$A$ (એટલે કે
$\Domain{M}$) પરનો પૂર્વક્રમ છે. બીજા શબ્દોમાં,~$\Gamma$ના નિદર્શો
બરાબર પૂર્વક્રમો છે. માત્ર~$\Obj P$ને !!{predicate} તરીકે ધરાવતી
પ્રથમ-ક્રમ ભાષામાં વ્યક્ત કરી શકાય એવો બધા પૂર્વક્રમોનો કોઈ પણ
ગુણધર્મ (ઉપરના સ્વવાચકતા અને પરંપરિતતાના ગુણધર્મોની જેમ)~$\Gamma$માંના
બે !!{sentence}sમાંથી ફલિત થાય છે અને તેનો ઊલટો પણ. એટલે
~$\Gamma$ના નિદર્શો વિશે આપણે જે કંઈ સાબિત કરીએ તે બધા પૂર્વક્રમો
વિશે સાબિત કર્યું છે.

દરેક નિશ્ચિત સિદ્ધાંત અને નિદર્શોના વર્ગ માટે (જેમ કે~$\Gamma$ અને
બધા પૂર્વક્રમો), સંબંધિત પ્રથમ-ક્રમ ભાષામાં શું વ્યક્ત કરી શકાય અને
શું ન કરી શકાય તે વિશે રસપ્રદ પ્રશ્નો હશે. !!{structure}sના કેટલાક
ગુણધર્મો બધી ભાષાઓ અને નિદર્શવર્ગો માટે રસપ્રદ છે; ખાસ કરીને
!!{domain}ના કદને લગતા ગુણધર્મો. દાખલા તરીકે, કોઈ પણ~$n \in \PosInt$
માટે !!{domain}માં બરાબર $n$~!!{element}s છે એવું હંમેશાં વ્યક્ત કરી
શકાય છે. અનંત સંખ્યામાં !!{sentence}sનો ગણ વાપરીને !!{domain} અનંત
છે એ પણ વ્યક્ત કરી શકાય છે. પરંતુ વ્યાપકક્ષેત્ર સાન્ત છે અથવા
વ્યાપકક્ષેત્ર !!{nonenumerable} છે એ વ્યક્ત કરી શકાતું નથી.
પ્રથમ-ક્રમ ભાષાઓની આ મર્યાદાઓ વિશેનાં પરિણામો સઘનતા અને
લેવેનહાઇમ--સ્કોલેમ પ્રમેયોનાં પરિણામો છે.

\end{document}
